% !TEX root = ../main.tex

\section{Kwantowe maszyny cieplne}
Kwantowe maszyny cieplne jako czynnik roboczy wykorzystują układ kwantowy o widmie ciągłym lub dyskretnym. Najprostszym przykładem jest układ o dwóch stanach energetycznych. Z układów o widmie dyskretnym powszechne są na przykład oscylator harmoniczny (energia stanu jest funkcją liniową liczby kwantowej) oraz cząstka w studni potencjały (energia stanu jest funkcją kwadratową liczby kwantowej).

Podobnie jak w przypadku klasycznych maszyn cieplnych zmiana energia może być zmieniana na dwa sposoby: ciepło i pracę. Praca jest dostarczana w sposów uporządkowany, a ciepło samorzutny. Według wzoru \eqref{eq:wartosc_oczekiwana_macierz_gestosci_mieszany} energia wewnętrzna układu opisanego hamiltonianem \(\ham\) i macierzą gęstości \(\vu{\rho}\) jest wartością oczekiwaną operatora Hamiltona:
\begin{equation*}
	U = \Tr\Bab{\vu{\rho}\ham}.
\end{equation*}
Różniczkując otrzymujemy:
\begin{equation*}
	\d{U} = \Tr\Bab{\d{\vu{\rho}}\ham} + \Tr\Bab{\vu{\rho}\d{\ham}}.
\end{equation*}
Z Pierwszej Zasady Termodynamiki mamy:
\begin{equation*}
	\d{U} = \idif{Q} + \idif{W}.
\end{equation*}
Hamiltonian \(\ham\) jest zależny od kontrolowanego przez nas parametru. Praca jest uporządkowana i możemy ją powiązać z członem zawierającym różniczkę hamiltonianu. Ciepło jest samoczynne i tu powiązane ze zmianą rozkładu opisanego macierzą gęstości \(\vu{\rho}\) \cite{myers}:
\begin{gather}
	\idif{Q} = \Tr\Bab{\d{\vu{\rho}}\ham}, \\
	\idif{W} = \Tr\Bab{\vu{\rho}\d{\ham}}.
\end{gather}

W stanie równowagi macierz gęstości opisana jest rozkładem Boltzmanna:
\begin{equation*}
	\vu{\rho}^{eq} = \inv{Z} \exp\pab{-\beta\ham}.
\end{equation*}